\section{OVERVIEW OF THE PROBLEM}
\label{ch:overview}

This chapter introduces credit-risk scoring on a peer-to-peer platform, explains why the problem is harder than it looks, and states the objectives and the approach followed in the rest of the report.

\subsection{Credit-risk scoring and peer-to-peer lending}

Credit scoring is the practice of turning information about an applicant into a number that orders applicants by their risk of not repaying; lenders use that number to accept or reject applications, to price loans and to set limits. The statistical literature on consumer credit scoring is long and dominated by logistic regression and scorecards because they are accurate enough, transparent to regulators and easy to monitor \cite{hand1997,siddiqi2006}. In peer-to-peer (P2P) lending, individual investors rather than a bank bear the credit risk, so the platform has to publish a risk signal: Lending Club assigns every loan a grade (A to G) and a sub-grade (A1 to G5, 35 levels) and ties the interest rate to it \cite{serrano2015}.

The public Lending Club file used here (Kaggle dataset \texttt{wordsforthewise/lending-club} \cite{kaggle_lc}) contains accepted loans issued between 2007 and 2018 with 151 columns: loan terms, borrower declarations, credit-bureau variables at the time of application, the repayment history after origination, hardship plans and debt-settlement fields. Two facts about the file shape the whole study. First, only some of the columns are knowable at application time (Section~\ref{sec:leakage}). Second, loans are observed until a snapshot date (December 2018 in this file), so recent loans have not had time to reach a final outcome (Section~\ref{sec:filters}).

\subsection{Challenges}

\paragraph{Data leakage.} A column is \emph{leaky} if it is generated by, or informative only because of, the outcome that we want to predict. In this dataset, columns such as \code{recoveries}, \code{total\_rec\_prncp} or \code{last\_fico\_range\_high} are recorded after origination. Section~\ref{sec:res-leak} measures how strong they are: in the training data the best post-origination column alone reaches a single-column ROC-AUC of 0.9074, versus 0.6565 for the best application-time column. A model that used them would look excellent and be useless for deciding on a new application.

\paragraph{Temporal dependence and dataset shift.} The relationship between borrower characteristics and default is not constant over calendar time: macro-economic conditions, platform policy and the mix of borrowers change. Mixing periods in a random split lets the model see the future when predicting the past, which inflates offline scores \cite{qian2021}. In our own data the default rate drifts from 13.70\% (train) to 14.74\% (validation) and 14.82\% (test), so even within three years the base rate moves.

\paragraph{Right-censoring and label maturity.} A loan issued in 2017 with a 60-month term cannot have a final outcome in a December 2018 snapshot. Treating ``not yet defaulted'' as ``repaid'' would bias the labels. Section~\ref{sec:filters} shows that, by issue year, the share of loans with a final label falls from 100\% (2011--2013) to 11.4\% (2018).

\paragraph{An incumbent model.} The platform's \code{sub\_grade} is itself a risk model built by the lender. It is both a highly informative application-time variable and a natural benchmark. A new model must be compared with it; putting it into a new model also changes how the other coefficients should be interpreted (Section~\ref{sec:res-exp}).

\paragraph{Messy and partly missing predictors.} Credit-bureau fields appear in the file at different dates, are partly missing even within the retained period (26 columns are 5--30\% missing in the training set), and are highly skewed (Section~\ref{sec:res-prep}); some pairs are almost duplicates (loan amount and instalment, correlation 0.994).

\paragraph{Moderate class imbalance.} After filtering, 13.7\% of training loans are defaults. This is imbalanced but not extreme; we therefore evaluate ranking quality (AUC, Gini, KS), which does not depend on a threshold or on the base rate.

\subsection{Motivation}

Beyond the practical interest of the problem, the dataset is a good teaching case for statistical learning. It requires the student to decide \emph{what may be used} (leakage), \emph{on whom the model is tested} (out-of-time evaluation), \emph{against what} (the incumbent), and \emph{how sure} one can be of an improvement (confidence intervals), rather than only to fit models. Many public projects on this dataset report high scores that mix periods or keep post-origination columns; the motivation here is to show what the numbers become when these shortcuts are removed.

\subsection{Objectives and approach}

The objective is to estimate, honestly, how much application-time information adds to the platform's own sub-grade for ranking default risk. The approach, detailed in Chapter~\ref{ch:method}, is:

\begin{enumerate}
    \item restrict to labelled, matured, 36-month loans issued from 2012 and split by issue date into train/validation/test (70/15/15);
    \item whitelist application-time columns using a documented dictionary, and verify the whitelist against single-column AUC;
    \item clean, engineer and encode features with a scikit-learn pipeline fitted only on the training set;
    \item compare a \code{sub\_grade} baseline with logistic regression, a decision tree, a random forest and histogram gradient boosting, tuned with time-ordered cross-validation inside the training set;
    \item quantify differences with paired-bootstrap intervals on the validation set and evaluate the final models once on the test set;
    \item study the role of \code{sub\_grade} (removal and coefficient shift) and the cost of a compact 12-field scorecard.
\end{enumerate}

\subsection{Chapter summary}

Credit scoring ranks applicants by default risk; on Lending Club data the main pitfalls are leaky post-origination columns, temporal drift, label censoring and the existence of the platform's own grade as an incumbent. The next chapter formalises the task, reviews related work and introduces the methods and metrics that are used.
